\section{Additional cosets at the penultimate step}
\label{acc-section}

The reference character used to construct the intermediate adelic
classes has no auxiliary character primes.  Consequently the only
additional datum at an intermediate step is a particular additive
coset, and this datum occurs only at the penultimate step.  We now
classify these cosets, prove their simultaneous existence on every
decorated lattice, and identify the resulting level subgroups.  This
proves the class comparison used at every step of the reference
construction, including its characters and stabilizer weights
\cite{alpo-uniform-source}.  The more general transfer between columns
with nontrivial auxiliary characters is a separate assertion.

\subsection{The selected rational primes are local uniformizers}

Write $\mathcal O=\mathcal O_E$, where $E$ is a CM field of degree
$2f$.  For a nonzero prime ideal $P$, put
$A_P=\mathcal O_{E,(P)}$ and $\kappa_P=\mathcal O/P$.
The valuation $v_P$ is normalized to have value one on a
uniformizer of the discrete valuation ring $A_P$.

\begin{lemma}\label{acc-uniformizer}
Let $p$ be a rational prime and let $\pi\in\mathcal O$ satisfy
\[
 \pi\overline\pi=p,\qquad (\pi)+(\overline\pi)=\mathcal O.
\]
Suppose that $f$ distinct prime ideals $P_1,\ldots,P_f$ divide
$(\overline\pi)$ and that $|\kappa_{P_i}|=p$ for every $i$.
Then
\[
 (\overline\pi)=\prod_{i=1}^f P_i,
 \qquad v_{P_i}(p)=1.
\]
In particular $p=|\kappa_{P_i}|$ is a uniformizer in $A_{P_i}$.
This assertion includes $p=2$.
\end{lemma}

\begin{proof}
Conjugation preserves the absolute norm of a principal ideal.
Taking ideal norms in $\pi\overline\pi=p$ therefore gives
\[
 N((\overline\pi))^2=p^{2f},
 \qquad N((\overline\pi))=p^f.
\]
The $P_i$ are distinct maximal ideals and hence pairwise coprime.
Thus $(\overline\pi)\subset\prod_iP_i$.
Both ideals have norm $p^f$, so the quotient of the larger ideal
by the smaller has one element.  They are equal.
Coprimality of $(\pi)$ and $(\overline\pi)$ gives
$v_{P_i}(\pi)=0$.  Taking valuations in the displayed product
identity now gives $v_{P_i}(p)=0+1=1$.
\end{proof}

The actual construction has two distinct rational primes $g_1,g_2$.
For each it supplies exactly the hypotheses of this lemma, with
$f$ selected prime ideals.  Let $\mathcal G$ be their disjoint union.
For $P\in\mathcal G$ write
\[
 t_P=|\kappa_P|\in\{g_1,g_2\}.
\]
We have proved $t_PA_P=PA_P$.  This equality is the precise
hypothesis needed below; membership $t_P\in P$ alone would not suffice.

\subsection{A quotient of intersections of local lattices}

The next statement applies to any number field, without a Hermitian
form.  Let $N\subset E^n$ be a full $\mathcal O$-lattice,
let $\mathcal G$ be a finite set of nonzero prime ideals, and
choose nonzero $t_P\in\mathcal O$ with $v_P(t_P)=1$.
Write $N_P=A_PN$, and fix a $\kappa_P$-linear isomorphism
\[
 b_P:N/PN\longrightarrow\kappa_P^n.
\]
It extends uniquely to $N_P/PN_P$ by inverting the reductions
of elements outside $P$.  We use $b_P$ also for this extended
reduction map on $N_P$.
Define $\mathcal O$-submodules of $E^n$ by
\[
 L(N)=\bigcap_{P\in\mathcal G}N_P,
 \qquad
 L^-(N)=\bigcap_{P\in\mathcal G}t_P^{-1}N_P.
\]
An intersection over an empty set is $E^n$.
These modules need not be finitely generated over $\mathcal O$.
Only their quotient, rather than either module separately, is finite.

\begin{theorem}\label{acc-residues}
The map
\begin{equation}\label{acc-reduction-map}
 L^-(N)\longrightarrow\mathcal A_n:=
       \prod_{P\in\mathcal G}\kappa_P^n,
 \qquad y\longmapsto \bigl(b_P(t_Py)\bigr)_P
\end{equation}
is surjective and has kernel $L(N)$.  It induces an isomorphism
of $\mathcal O$-modules
\[
                    L^-(N)/L(N)\simeq\mathcal A_n.
\]
In particular every element of $\mathcal A_n$ is realized on every
lattice $N$, and
\[
 |L^-(N)/L(N)|=\prod_{P\in\mathcal G}|\kappa_P|^n.
\]
\end{theorem}

\begin{proof}
Since $t_PA_P=PA_P$, the equality $b_P(t_Py)=0$ is equivalent to
$t_Py\in t_PN_P$, and hence to $y\in N_P$.
This proves the kernel assertion and also shows that
$L(N)\subset L^-(N)$.

For surjectivity, prescribe vectors $a_P\in\kappa_P^n$.
Choose $x_P\in N$ with $b_P(x_P)=a_P$.
For each fixed $P$, the Chinese remainder theorem gives
$e_P\in\mathcal O$ such that
\[
 e_P\equiv1\pmod P,\qquad
 e_P\in Q^{v_Q(t_P)}\quad
       (Q\in\mathcal G,\ Q\ne P).
\]
An exponent zero imposes no condition.  The moduli in these
conditions are pairwise coprime, so the theorem applies.
Put
\[
                         y=\sum_{P\in\mathcal G}
                                      \frac{e_P}{t_P}x_P.
\]
Fix $Q\in\mathcal G$.  The $P=Q$ term in $t_Qy$ is
$e_Qx_Q\in N_Q$ and has reduction $a_Q$.
For $P\ne Q$ the chosen divisibility gives
$e_P/t_P\in A_Q$, so
\[
             t_Q(e_P/t_P)x_P\in t_QN_Q=QN_Q.
\]
Thus $t_Qy\in N_Q$ and $b_Q(t_Qy)=a_Q$.
This proves surjectivity simultaneously at all the selected primes.
When $\mathcal G$ is empty, both the quotient and the product have
one element and $y=0$ suffices.  The same proof permits $n=0$.
The first isomorphism theorem and the cardinalities of the residue
vector spaces prove the remaining assertions.
\end{proof}

For the actual two rational primes the number in this theorem is
\begin{equation}\label{acc-number-cosets}
                         |\mathcal A_n|=(g_1g_2)^{fn}.
\end{equation}
A permitted additional coset is
\[
                          Y=y+L(N),\qquad y\in L^-(N).
\]
Its residue vector, denoted $a(Y)$, is the vector in
\eqref{acc-reduction-map}.  It does not depend on $y\in Y$.
The theorem says precisely that $Y\mapsto a(Y)$ is a bijection.

\subsection{Coordinates and change of residue frames}

Choose an $A_P$-basis matrix $F_P$ of $N_P$ that lifts the
inverse of $b_P$; explicitly,
\[
 F_P A_P^n=N_P,
 \qquad b_P(F_Px)=\overline x\quad(x\in A_P^n).
\]
Such a matrix exists: choose any local basis, then correct its
residue matrix by a lift in $\operatorname{GL}_n(A_P)$.
The reduction map on this general linear group is onto, since
any entrywise lift of an invertible residue matrix has unit determinant.
Define the coset in these coordinates by
\begin{equation}\label{acc-coordinate-coset}
                    C_P(Y)=F_P^{-1}y+A_P^n\subset E^n.
\end{equation}

\begin{proposition}\label{acc-coordinate-classification}
The coset \eqref{acc-coordinate-coset} is independent of the
representative $y\in Y$ and of the compatible basis matrix $F_P$.
For permitted cosets on two possibly different lattices with
residue frames,
\[
 \bigl(C_P(Y)\bigr)_{P\in\mathcal G}
   =\bigl(C_P(Y')\bigr)_{P\in\mathcal G}
                   \quad\Longleftrightarrow\quad a(Y)=a(Y').
\]
\end{proposition}

\begin{proof}
Changing $y$ by an element of $L(N)$ changes $F_P^{-1}y$ by
an element of $A_P^n$.  Moreover
$t_PF_P^{-1}y\in A_P^n$ and its reduction is $a_P(Y)$.
For $u,u'\in t_P^{-1}A_P^n$,
\[
 u+A_P^n=u'+A_P^n
   \ \Longleftrightarrow\ u-u'\in A_P^n
   \ \Longleftrightarrow\ \overline{t_Pu}=\overline{t_Pu'}.
\]
The last equivalence uses $t_PA_P=PA_P$.
This already proves the asserted comparison between two lattices.
For completeness, another compatible basis is $F'_P=F_Pk$,
where $k\in\operatorname{GL}_n(A_P)$ has reduction one.
Then $k^{-1}-1\in t_PM_n(A_P)$; hence
$(k^{-1}-1)F_P^{-1}y\in A_P^n$.
The coordinate coset is unchanged.
\end{proof}

Return to the positive Hermitian space and decorated lattices of
Section~\ref{acl-section}.  Here the selected residue-frame primes
are $\mathcal G$ and the functional primes are disjoint from them.
Let $\mathscr D$ be the set of decorated lattices and let
\[
                   \mathscr P=\{(d,Y):d\in\mathscr D,
                        \ Y\in L^-(N_d)/L(N_d)\}.
\]
The notation treats a quotient element as its additive coset.
The rational unitary group and the residue general linear groups act by
\[
 (\gamma,h)(d,Y)=((\gamma,h)d,\gamma Y),
 \qquad h=(h_P)_P\in\prod_P\operatorname{GL}_n(\kappa_P).
\]
Recall that the new residue frame is
$b'_P=h_Pb_P\gamma^{-1}$.
Since $t_P\in E$, it commutes with the $E$-linear map $\gamma$.
Consequently
\begin{equation}\label{acc-equivariance}
                a_P(\gamma Y)=h_Pa_P(Y).
\end{equation}
Theorems~\ref{acc-residues} and \ref{acc-coordinate-classification}
therefore give a bijection, compatible with both actions,
\begin{equation}\label{acc-product-description}
 \mathscr P\simeq\mathscr D\times\mathcal A_n,
 \qquad
 (\gamma,h)(d,a)=((\gamma,h)d,ha).
\end{equation}
In particular a rational unitary change alone preserves the coordinate
cosets, even when the chosen lifts of the residue frames are different.
A residue-frame change acts on their labels by the ordinary matrix
action.  No choice of global basis of $N_d$ is involved.

\subsection{The exact partition into adelic classes}

Let $B$ be a finite set of indices.  For $b\in B$ let
\[
 H_b=\prod_{P\in\mathcal G}H_{b,P}
      \subset\prod_{P\in\mathcal G}\operatorname{GL}_n(\kappa_P)
\]
be a product of subgroups, and let $\chi_b:H_b\to K^\times$
be a character in the coefficient field.  For $a\in\mathcal A_n$
put
\begin{equation}\label{acc-stabilizer-level}
 H_{b,a}=\{h\in H_b:ha=a\}
      =\prod_{P\in\mathcal G}
            \bigl(H_{b,P}\cap\operatorname{Stab}(a_P)\bigr).
\end{equation}
Here $\operatorname{Stab}(a_P)$ fixes the vector $a_P$, not merely
its one-dimensional span.  Choose one representative $a$ of each
$H_b$-orbit in $\mathcal A_n$.

\begin{theorem}\label{acc-class-partition}
There is a bijection
\begin{equation}\label{acc-class-bijection}
 \coprod_{b\in B}\ \coprod_{a\in H_b\backslash\mathcal A_n}
       (U(F)\times H_{b,a})\backslash\mathscr D
 \ \longrightarrow\
 \coprod_{b\in B}(U(F)\times H_b)\backslash\mathscr P,
 \qquad [d]\longmapsto[(d,a)].
\end{equation}
Every stabilizer on the left is isomorphic to the stabilizer of
its corresponding point on the right.  These isomorphisms preserve
the rational unitary component and the character, whose restriction
on the left is $\chi_b|_{H_{b,a}}$.

All class sets in \eqref{acc-class-bijection} are finite.
The number of displayed pieces is at most
$|B|\,|\mathcal A_n|$.
\end{theorem}

\begin{proof}
Use the product description \eqref{acc-product-description}.
Given $(d,a')$, choose $h\in H_b$ carrying $a'$ to the selected
representative $a$.  Applying $(1,h)$ puts its class in the indicated
piece.  Distinct $H_b$-orbits of vectors cannot meet.
If $(\gamma,h)(d,a)=(d',a)$, then $ha=a$, so
$h\in H_{b,a}$; this proves injectivity within the piece.
The same observation identifies the stabilizer groups themselves:
\[
 \operatorname{Stab}_{U(F)\times H_b}(d,a)
       =\operatorname{Stab}_{U(F)\times H_{b,a}}(d).
\]
It does not only identify their orders or their images in $U(F)$.
Both the rational component and the restricted character are unchanged.

Each group $H_{b,a}$ is a finite residue subgroup, so the class
finiteness theorem of Section~\ref{acl-section} applies.
There are finitely many $b$ and finitely many vector orbits.
There are at most $|\mathcal A_n|$ vector orbits for each $b$,
which proves the quantitative assertion.
\end{proof}

At every step other than the penultimate step the additional datum is
the empty set, rather than a nonempty additive coset.
For those steps replace $\mathcal A_n$ in the theorem by a one-element
set with trivial action.  Then $H_{b,a}=H_b$ and there is exactly
one piece for each $b$.  In the actual construction $m\geq2$,
so both the initial and the final step are of this kind.
Their index set $B$ has one element, giving a unique piece at each end.
At the penultimate step the number of pieces is bounded by
\[
                            |B|(g_1g_2)^{fn}.
\]

The source chooses representatives by first fixing the cosets in local
coordinates.  Proposition~\ref{acc-coordinate-classification} shows
that this is equivalent to fixing the vectors $a_P$ above.
The orbit of such a vector depends only on its $H_b$-orbit and not on
the underlying decorated lattice, by the simultaneous surjectivity
in Theorem~\ref{acc-residues}.  Thus the choice of preferred
representatives introduces no dependence on an unavailable lattice
or later parameter.  It selects exactly the pieces in
\eqref{acc-class-bijection}, up to renaming their indices.

\subsection{Level subgroups, coefficients, and all volume factors}

Use the notation of Section~\ref{acl-section}:
$K_S$ is the compact open subgroup preserving the base lattice and
its fixed residue functionals, and
\[
 r:K_S\longrightarrow
           \prod_{P\in\mathcal G}\operatorname{GL}_n(\kappa_P)
\]
is the surjective residue map.  Put
$K_{b}=r^{-1}(H_b)$ and $K_{b,a}=r^{-1}(H_{b,a})$.
Theorem~\ref{acl-double-cosets} identifies the $(b,a)$ piece with
\begin{equation}\label{acc-adelic-piece}
                         U(F)\backslash U(\mathbb A_{F,f})/K_{b,a}.
\end{equation}
The exact level character is
\begin{equation}\label{acc-character}
       \omega_{b,a}(k)=\chi_b(r(k))^{-1}
                              \qquad(k\in K_{b,a}).
\end{equation}
Indeed right multiplication by $k$ changes the residue frames by
$r(k)^{-1}$, so this follows directly from
\eqref{acc-equivariance} and the equivariance of coefficient functions.
No additional character comes from the additive coset.

More explicitly, suppose coefficient functions on $\mathscr P$ satisfy
\[
 F((\gamma,h)(d,a))=\chi_b(h)R(\gamma)F(d,a),
\]
where $R$ is the polynomial representation in
Section~\ref{acl-section}.  Restriction to a fixed representative $a$
gives exactly the corresponding equivariant coefficient function on
$\mathscr D$ for $U(F)\times H_{b,a}$.
Conversely, if $F_a$ is the function on the chosen piece, choose
$g=(\gamma,h)$ with $ha=a'$ and set
\[
 F(d,a')=\chi_b(h)R(\gamma)F_a(g^{-1}d).
\]
Two choices differ by an element of $U(F)\times H_{b,a}$,
so equivariance makes the result independent of the choice.
This proves the coefficient-space decomposition.

The integral lattice attached to a point depends on its functional
kernel lattice $L_d$ and on its fixed weight.  It does not depend on
$Y$.  Therefore the same restriction and transport identify the
integral coefficient lattices, as well as their rational and complex
spans.  In particular, the conjugate-prime coefficient identity of
Section~\ref{acl-section} remains valid on every piece.
If stabilizers are nontrivial their invariant subspaces are retained.
The weighted Hermitian norm is the orthogonal sum of the piece norms,
because corresponding stabilizers are the identical groups in
Theorem~\ref{acc-class-partition}.

\begin{proposition}\label{acc-index-volume}
For each vector $a$,
\[
 [K_b:K_{b,a}]=[H_b:H_{b,a}]=|H_ba|.
\]
For any one fixed Haar measure $\mu$ on the finite adelic group,
\[
                     \mu(K_{b,a})=\frac{\mu(K_b)}{|H_ba|}.
\]
The quotient integral for the piece
\eqref{acc-adelic-piece} is consequently
$\mu(K_{b,a})$ times its finite class sum with weights
$1/|\operatorname{Stab}|$.
\end{proposition}

\begin{proof}
Surjectivity of $r$ identifies the coset spaces
$K_b/K_{b,a}$ and $H_b/H_{b,a}$.
The map $hH_{b,a}\mapsto ha$ is a bijection from the latter
coset space to the orbit of $a$.
The finite disjoint union of $[K_b:K_{b,a}]$ translates of
$K_{b,a}$ has measure $\mu(K_b)$, giving the volume formula.
The quotient integral formula is the one proved in
Section~\ref{acl-section}, now at the subgroup $K_{b,a}$.
\end{proof}

This proof also applies to the source's local coset stabilizer:
let $A$ be the ring of integers of a nonarchimedean local field $k$,
with uniformizer $\varpi$.  If $J\subset\operatorname{GL}_n(A)$
preserves a lattice $\Lambda\subset k^n$, it acts on $k^n/\Lambda$; the index in $J$ of the
stabilizer of a quotient element is exactly its orbit size whenever
that orbit is finite.  One proves this by the same bijection from
stabilizer cosets to orbit elements.  For $\Lambda=\varpi^M A^n$, with $M\geq0$ an integer, and
vectors in $\varpi^{-M}A^n$, the quotient containing the orbit is finite.
Thus the relative-index factor in the additional local-level weight
is an orbit cardinality, with no omitted inverse.

There is a useful direct check on the stabilizer weights.
For a fixed decorated class represented by $d$, let
$\Gamma=\operatorname{Stab}_{U(F)\times H_b}(d)$.
Its inverse image in the enhanced class set consists of the
$\Gamma$-orbits in $\mathcal A_n$.
For any function $w$ constant on those orbits, orbit--stabilizer gives
\begin{equation}\label{acc-weighted-fiber}
 \sum_{[a]\in\Gamma\backslash\mathcal A_n}
           \frac{w(a)}{|\Gamma_a|}
      =\frac1{|\Gamma|}\sum_{a\in\mathcal A_n}w(a).
\end{equation}
To prove it, expand the sum on the right as a sum over orbits;
each orbit has $|\Gamma|/|\Gamma_a|$ elements.
In particular adjoining all the additional cosets multiplies the
weighted mass of a decorated class by $|\mathcal A_n|$.
Restricting to one piece instead uses its actual stabilizer and does
not insert this factor.  These two different operations must not be
confused.

\subsection{Checks on the necessary hypotheses}

\begin{example}\label{acc-norm-counterexample}
The induced map $L^-(N)/L(N)\to\mathcal A_n$ in
Theorem~\ref{acc-residues} need not be injective if
$v_P(t_P)=1$ is replaced by $t_P\in P$.
This can fail even when the number field is unramified at the prime.
Take $E=\mathbb Q(i)$, $P=(3)$, $A=A_P$ and $N_P=A$.
The polynomial $X^2+1$ is irreducible modulo $3$, so
$|\kappa_P|=9$ and $3$ is a uniformizer in $A$.
Set $t_P=|\kappa_P|=9$.
The cosets $A$ and $1/3+A$ are distinct, because $1/3\notin A$,
but both have residue $\overline{9y}=0$.
Thus the size of the residue field cannot be used as a uniformizer
at an arbitrary prime.  Lemma~\ref{acc-uniformizer} proves exactly
the stronger condition used at the actual selected primes.
\end{example}

\begin{example}\label{acc-ramified-counterexample}
A residue-degree-one hypothesis by itself also does not suffice.
In $E=\mathbb Q(\sqrt2)$ at $P=(\sqrt2)$, the residue field is
$\mathbb F_2$, but $2=(\sqrt2)^2$ has valuation two.
The distinct cosets $A_P$ and $1/\sqrt2+A_P$ both map to zero
under $y\mapsto\overline{2y}$.
The coprime factorization in Lemma~\ref{acc-uniformizer}
excludes this ramification at the selected primes.
\end{example}

\begin{example}\label{acc-line-counterexample}
Replacing the vector stabilizer by a line stabilizer changes the
level and its index.  On $\mathbb F_3$, let
$H=\mathbb F_3^\times$ act by multiplication and let $a=1$.
Its vector stabilizer is trivial and its orbit has two elements.
Its line stabilizer is all of $H$, since the vector space has
dimension one.  The resulting indices are two and one respectively.
The actual coordinate datum fixes an additive coset and hence a
vector, so it requires the former subgroup.
\end{example}

\begin{example}\label{acc-weight-counterexample}
In the same action of $\Gamma=\{\pm1\}$ on $\mathbb F_3$,
there are two orbits, $\{0\}$ and $\{1,-1\}$.
Their stabilizer orders are two and one.  Their weighted sum is
$1/2+1=3/2$, as \eqref{acc-weighted-fiber} predicts.
Using the original stabilizer order for both orbits would give one.
Thus an unchanged stabilizer denominator after adjoining a coset
would be incorrect; the actual chart comparison uses the smaller
stabilizer and has the stated exact equality.
\end{example}

These examples concern weakened statements, not flaws in the
actual source hypotheses.  The argument above proves the additional
coset classification, its simultaneous realization, the finite
partition into reference-step classes, and their exact integral
coefficients, characters, stabilizers and volume factors.
It uses the selected-prime hypotheses already verified in the
arithmetic construction.  It does not assert a comparison for
arbitrary nontrivial auxiliary character primes, nor does it prove
the remaining theta-period identities or the automorphic spectral
and mass formulas.


